\section{Renewal bounds and complete cancellation}
Let \(U=\mathbf1_{[0,1)}\) and write
\(\delta_0+\sum_{m\ge1}U^{*m}=\delta_0+r(t)dt\).
The renewal series converges locally by
\(U^{*m}(t)\le t^{m-1}/(m-1)!\). Right values are used at jumps.
The renewal equation gives
\[
 r(t)=e^t\ (0\le t<1),\qquad
 r(t)=e^t-te^{t-1}\ (1\le t<2),\qquad
 r(t)=\int_{t-1}^t r(s)ds\ (t\ge1).
\]
Define \(b=(r/2-1)\mathbf1_{[0,\infty)}\) and
\(\sigma=(r/2)\mathbf1_{[0,\infty)}\).

\begin{lemma}\label{renewal}
The following bounds hold:
\begin{gather*}
 1\le r\le e,\quad \|b\|_\infty\le\tfrac12,\quad
 \|b\|_1\le4,\quad \TV(b)\le10,\quad \int b=-\tfrac16,\\
 |b(t)|\le\tfrac43e^{-t/4},\qquad
 \int_v^\infty|b(t)|dt\le6e^{-v/4}.
\end{gather*}
Writing \(A_b(s)=\int_0^\infty b(t)b(t+s)dt\), one has, for \(s>0\),
\begin{equation}
 \mathcal B(s):=\frac16-\int_0^s b+\frac12b(s)+A_b(s)\le\frac s2.
 \label{cancel}
\end{equation}
For each real \(U\), the function
\begin{equation}
 C_U(s)=U-s-\frac23-2\int_0^s b+2A_b(s)
 \label{CU}
\end{equation}
is 12-Lipschitz on \(s>0\), and \(C_U(s)+\sigma(s)\le U\).
\end{lemma}
\begin{proof}
Here are explicit reconstructions of the renewal estimates and Akwei's
cancellation inequality. For \(f(z)=r(m-z)\), the Volterra equation on the
next unit block gives
\[
 r(m+y)=\int_0^1q_y(z)f(z)dz,\quad
 q_y(z)=\mathbf1_{y+z\le1}
 +\int_{\max(0,y+z-1)}^y e^s ds.
\]
The kernel has mass one and is at least \(1/2\) on \(0\le z\le1/2\).
Thus the closed block ranges stay in \([1,e]\), are nested, and their
oscillations contract by \(3/4\). The invariant
\(\int_{t-1}^t(s-t+1)r(s)ds=1\) identifies the limit as 2. Hence
\(|r(t)-2|\le(e-1)(3/4)^{\lfloor t\rfloor}\).
This proves the stated norm and tail bounds, using \(e<3\) and
\(\log(4/3)>1/4\). The jumps of \(b\) at 0 and 1 contribute 1 to its
variation. The regular variation on \((0,1)\) is less than 1, and
\(b'=b(t)-b(t-1)\) on \(t>1\) contributes at most \(2\|b\|_1\le8\).

Put \(g=2b\), \(G(t)=\int_0^t g\), and
\(A_g(s)=\int_0^\infty g(t)g(t+s)dt\).
The Laplace transform
\[
 \widehat r(z)=\frac{1-e^{-z}}{z-1+e^{-z}}
 =\frac2z-\frac13+O(z)
\]
and exponential integrability give \(\int g=-1/3\).
Splitting at the moving jump before differentiation gives, for \(0<s<1\),
\[
 A_g'(s)=A_g(s)-A_g(1-s)+e^{1-s}+4s-4.
\]
Subtracting at \(1-s\) and using antisymmetry gives
\(A_g(s)-A_g(1-s)=e^s-e^{1-s}-4s+2\). Thus
\(A_g(s)=e^s-2s+C\). On \(s>1\), \(A_g'=A_g(s)-A_g(s-1)\).
Decay forces both its moving-integral invariants to vanish, in particular
\(\int_0^1 sA_g(s)ds=0\). Consequently \(C=-2/3\) and \(A_g(0)=1/3\).

The elementary renewal stopping-time identity gives \(G(t)\ge-1\): if
\(T\) is the first sum of independent uniform increments exceeding \(t\),
then \(E T<\infty\), \(E\sum_{j\le T}U_j=E T/2>t\), and
\(\int_0^t r=E T-1\).
Set \(h(s)=2s+2+2G(s)-g(s)-A_g(s)\). The formulas above yield \(h=8/3\)
on \(0<s<1\). For \(s\ge1\),
\(h=2s+2+G(s)+G(s-1)-A_g(s)\), with \(A_g(s)\le1/3\).
For \(s\ge3/2\) this proves \(h\ge8/3\). For \(1\le s\le3/2\),
\(G(s)\ge3(e-3)/2\), \(G(s-1)\ge\sqrt e-2\), so
\(h\ge3e/2+\sqrt e-17/6>8/3\).
These are exactly \eqref{cancel}, including the right value at 1.
Finally BV translation gives \(\operatorname{Lip}(A_b)\le5\), so
\eqref{CU} is 12-Lipschitz and
\(C_U+\sigma=U-s+2\mathcal B\le U\).
\end{proof}

The purpose of retaining the whole last inequality is that \(\sigma\) jumps
at 1. The proof below translates \(C_U\), which is Lipschitz, rather than
Taylor expanding the discontinuous factor.
